\documentclass[aspectratio=169]{beamer}
\usetheme{Madrid}
\usecolortheme{default}
\usepackage[utf8]{inputenc}
\usepackage[T5]{fontenc}
\usepackage{graphicx}
\usepackage{hyperref}
\usepackage{booktabs}
\usepackage{amsmath}

\setbeamertemplate{caption}[numbered]
\renewcommand{\figurename}{Hình}
\renewcommand{\thefigure}{12.\arabic{figure}}

\title[Chương 12: Object Detection]{Học Sâu (Deep Learning) \\ \vspace{0.5cm} \Large Chương 12: Phát hiện Vật thể \\ (Object Detection)}
\author{Đại học Đông Á}
\date{\today}

\begin{document}

\begin{frame}
    \titlepage
\end{frame}

\begin{frame}{Nội dung chương}
    \tableofcontents
\end{frame}

\section{1. Giới thiệu Bài toán Phát hiện Vật thể}

\begin{frame}{Bài toán Phát hiện Vật thể là gì?}
    \begin{columns}
        \column{0.5\textwidth}
        \begin{itemize}
            \item Không chỉ trả lời câu hỏi "Trong ảnh có gì?" (Classification).
            \item Phải trả lời được "Nó nằm ở đâu?" bằng cách vẽ các \textbf{Hộp giới hạn (Bounding Boxes)}.
            \item \textbf{Ứng dụng thực tế:}
            \begin{itemize}
                \item Xe tự lái: Phát hiện người đi bộ, biển báo.
                \item Bán lẻ: Đếm số lượng sản phẩm trên kệ.
                \item Trích xuất: Cắt vùng biển số xe để đưa vào OCR.
            \end{itemize}
        \end{itemize}
        
        \column{0.5\textwidth}
        \begin{figure}
            \centering
            \includegraphics[width=\textwidth,height=0.6\textheight,keepaspectratio]{../../images/ch12/object-detection.7c5cbfd4.png}
            \caption{Vẽ hộp giới hạn và gán nhãn}
        \end{figure}
    \end{columns}
\end{frame}

\begin{frame}{Bộ dữ liệu COCO (Common Objects in Context)}
    \begin{columns}
        \column{0.5\textwidth}
        \begin{itemize}
            \item MS COCO là tiêu chuẩn vàng (Gold Standard) để đánh giá các mô hình Object Detection hiện đại.
            \item Hơn 200,000 hình ảnh được gán nhãn thủ công với 80 danh mục vật thể.
            \item Tính chất đa dạng cao: Nhiều vật thể đan xen, chồng chéo, kích thước từ siêu nhỏ đến khổng lồ trong cùng một khung hình.
        \end{itemize}
        
        \column{0.5\textwidth}
        \begin{figure}
            \centering
            \includegraphics[width=\textwidth,height=0.6\textheight,keepaspectratio]{../../images/ch12/coco-example.c43c39d6.png}
            \caption{Dữ liệu phức tạp trong COCO}
        \end{figure}
    \end{columns}
\end{frame}

\section{2. Kiến trúc Hai giai đoạn (Two-stage Detectors)}

\begin{frame}{Họ mô hình R-CNN (Region-based CNN)}
    \begin{columns}
        \column{0.5\textwidth}
        Kiến trúc Hai giai đoạn ưu tiên \textbf{Độ chính xác cao}:
        \begin{enumerate}
            \item \textbf{Giai đoạn 1 (Đề xuất vùng - Region Proposals):} Quét toàn bộ ảnh để tìm ra vài ngàn Vùng (Regions) "có vẻ chứa vật thể".
            \item \textbf{Giai đoạn 2 (Phân loại):} Đưa từng vùng đó qua mạng CNN để dự đoán Nhãn (Label) và Tinh chỉnh tọa độ hộp (Box Refinement).
        \end{enumerate}
        \textit{Nhược điểm: Rất chậm, không thể chạy Real-time.}
        
        \column{0.5\textwidth}
        \begin{figure}
            \centering
            \includegraphics[width=\textwidth,height=0.6\textheight,keepaspectratio]{../../images/ch12/r-cnn-pipeline.8fe83666.png}
            \caption{Quy trình xử lý của R-CNN}
        \end{figure}
    \end{columns}
\end{frame}

\section{3. Kiến trúc Một giai đoạn (Single-stage Detectors)}

\begin{frame}{Kiến trúc Single-stage và FPN}
    \begin{columns}
        \column{0.5\textwidth}
        \begin{itemize}
            \item Dự đoán song song Tọa độ hộp và Nhãn ngay trong 1 bước duy nhất (YOLO, RetinaNet, SSD).
            \item \textbf{Vấn đề:} Vật thể nhỏ dễ bị mất dấu ở các lớp sâu.
            \item \textbf{Giải pháp FPN (Feature Pyramid Network):} Hợp nhất đặc trưng ở nhiều độ phân giải khác nhau, giúp mạng phát hiện được vật thể ở mọi kích cỡ (Từ nhỏ xíu đến to lớn).
        \end{itemize}
        
        \column{0.5\textwidth}
        \begin{figure}
            \centering
            \includegraphics[width=\textwidth,height=0.6\textheight,keepaspectratio]{../../images/ch12/feature-pyramid-network.83b6f108.png}
            \caption{Mạng tháp đặc trưng (FPN)}
        \end{figure}
    \end{columns}
\end{frame}

\begin{frame}{Kết quả dự đoán của RetinaNet}
    \begin{columns}
        \column{0.5\textwidth}
        \begin{itemize}
            \item RetinaNet là một trong những Single-stage Detector xuất sắc nhất của Facebook AI.
            \item Sử dụng hàm \textbf{Focal Loss} để giải quyết bài toán mất cân bằng lớp trầm trọng giữa Tiền cảnh (Chứa vật thể) và Hậu cảnh (Không chứa gì).
            \item Dự đoán mượt mà, bao quát toàn bộ khung hình một cách chính xác.
        \end{itemize}
        
        \column{0.5\textwidth}
        \begin{figure}
            \centering
            \includegraphics[width=\textwidth,height=0.6\textheight,keepaspectratio]{../../images/ch12/retinanet-output.0a67b6e8.png}
            \caption{Dự đoán của RetinaNet}
        \end{figure}
    \end{columns}
\end{frame}

\section{4. Thuật toán YOLO (You Only Look Once)}

\begin{frame}{Kỷ nguyên mới mang tên YOLO}
    \begin{itemize}
        \item Được giới thiệu năm 2015 bởi Joseph Redmon, YOLO làm rung chuyển giới AI.
        \item Chuyển bài toán Detection thành \textbf{Regression (Hồi quy) trên lưới (Grid)}.
        \item Khả năng chạy Real-time lên tới 45-155 FPS trên GPU, vượt xa hoàn toàn R-CNN.
        \item Nguyên lý: Chia ảnh thành lưới $S \times S$ (VD: $13 \times 13$). Mỗi ô lưới (Cell) chịu trách nhiệm dự đoán các vật thể có tâm (Center) nằm trong ô đó.
    \end{itemize}
\end{frame}

\begin{frame}{Tư duy Chia lưới và Anchor Boxes}
    \begin{columns}
        \column{0.5\textwidth}
        \begin{itemize}
            \item Hình bên trái: Chia ảnh gốc thành Lưới (Grid). Ô màu cam chứa Tâm của 1 con chó. Do đó ô lưới màu cam phải chịu trách nhiệm dự đoán con chó này.
            \item Hình bên phải: Mô hình sẽ không dự đoán tọa độ ngẫu nhiên. Nó sẽ chỉnh sửa dựa trên các \textbf{Hộp neo (Anchor Boxes)} có sẵn (Tỷ lệ chiều cao / chiều rộng).
        \end{itemize}
        
        \column{0.5\textwidth}
        \begin{figure}
            \centering
            \includegraphics[width=\textwidth,height=0.5\textheight,keepaspectratio]{../../images/ch12/yolo-diagram.b7347ca6.png}
            \caption{Cơ chế hoạt động của YOLO}
        \end{figure}
    \end{columns}
\end{frame}

\begin{frame}{Toán học: Cấu trúc Vector Đầu ra}
    \begin{columns}
        \column{0.5\textwidth}
        Mỗi ô lưới xuất ra một Vector gồm:
        \begin{enumerate}
            \item \textbf{$p_c$}: Confidence Score (Xác suất có chứa vật thể).
            \item \textbf{$b_x, b_y$}: Tọa độ tâm của hộp giới hạn.
            \item \textbf{$b_w, b_h$}: Chiều rộng, chiều cao của hộp.
            \item \textbf{$c_1, c_2, ... c_k$}: Xác suất thuộc về class $i$ (Vd: Chó, Mèo, Xe).
        \end{enumerate}
        Toàn bộ được huấn luyện song song trong 1 hàm Loss duy nhất.
        
        \column{0.5\textwidth}
        \begin{figure}
            \centering
            \includegraphics[width=\textwidth,height=0.6\textheight,keepaspectratio]{../../images/ch12/yolo-targets.7b3c7aed.png}
            \caption{Vector mục tiêu huấn luyện}
        \end{figure}
    \end{columns}
\end{frame}

\begin{frame}{Kết quả Thô (Raw Predictions)}
    \begin{columns}
        \column{0.5\textwidth}
        \begin{itemize}
            \item Trong quá trình dự đoán (Inference), MỖI ô lưới sẽ xuất ra nhiều Bounding Boxes.
            \item Đối với một ảnh $416 \times 416$, mô hình có thể sinh ra hơn \textbf{10,000 hộp giới hạn} cùng một lúc.
            \item Hầu hết các hộp này là nhiễu, chồng chéo lên nhau hoặc chỉ là nền (Background).
            \item Làm sao để dọn dẹp đống hỗn độn này?
        \end{itemize}
        
        \column{0.5\textwidth}
        \begin{figure}
            \centering
            \includegraphics[width=\textwidth,height=0.6\textheight,keepaspectratio]{../../images/ch12/yolo-predictions-all.de18d520.png}
            \caption{Hàng ngàn hộp dự đoán thô}
        \end{figure}
    \end{columns}
\end{frame}

\begin{frame}{Công cụ 1: Intersection over Union (IoU)}
    \begin{itemize}
        \item \textbf{IoU} là công thức toán học cốt lõi để tính toán độ trùng khớp giữa 2 Hộp giới hạn.
        \item IoU = (Diện tích phần Giao nhau) / (Diện tích phần Hợp nhau).
        $$ IoU = \frac{Area(Box_1 \cap Box_2)}{Area(Box_1 \cup Box_2)} $$
        \item IoU dao động từ $0$ (Không chạm nhau) đến $1$ (Trùng khớp hoàn toàn).
        \item Ngưỡng tiêu chuẩn: Thường tính là "Dự đoán đúng" nếu $IoU > 0.5$.
    \end{itemize}
\end{frame}

\begin{frame}{Công cụ 2: Non-Maximum Suppression (NMS)}
    Thuật toán NMS để lọc nhiễu hộp giới hạn:
    \begin{enumerate}
        \item Bỏ tất cả các hộp có Điểm tự tin (Confidence Score) $< 0.5$.
        \item \textbf{Vòng lặp:} Chọn hộp có Confidence Score \textbf{Cao nhất} làm chuẩn (Vd: Hộp chó $0.98$).
        \item Tính \textbf{IoU} của hộp chuẩn với tất cả các hộp Chó khác còn lại.
        \item Xóa toàn bộ các hộp có $IoU > 0.4$ (Vì chúng chắc chắn là bản sao trùng lặp của hộp chuẩn).
        \item Lặp lại đến khi không còn hộp nào chưa xét.
    \end{enumerate}
\end{frame}

\begin{frame}{Kết quả Cuối cùng (Sau khi chạy NMS)}
    \begin{columns}
        \column{0.5\textwidth}
        \begin{itemize}
            \item Sau khi trải qua bộ lọc Non-Maximum Suppression, từ hàng vạn hộp nhiễu, ta chỉ còn lại chính xác số lượng vật thể trong hình.
            \item YOLO vừa nhanh, vừa chính xác, thiết lập tiêu chuẩn công nghiệp (Industry Standard) cho mọi ứng dụng Camera AI, Flycam, và Xe tự lái hiện nay.
        \end{itemize}
        
        \column{0.5\textwidth}
        \begin{figure}
            \centering
            \includegraphics[width=\textwidth,height=0.6\textheight,keepaspectratio]{../../images/ch12/yolo-predictions.d621592d.png}
            \caption{Kết quả sạch sẽ sau NMS}
        \end{figure}
    \end{columns}
\end{frame}

\section{Tổng kết}

\begin{frame}{Tổng kết Chương 12}
    \begin{itemize}
        \item \textbf{Object Detection} phức tạp hơn Classification do phải giải quyết bài toán định vị (Localization).
        \item \textbf{Two-stage (R-CNN):} Quét vùng trước, Phân loại sau. Chậm nhưng chính xác.
        \item \textbf{Single-stage (YOLO, RetinaNet):} Gộp chung thành bài toán Hồi quy (Regression) trên lưới tọa độ. Nhanh và hiệu quả.
        \item Hiểu rõ các nguyên lý Toán học ẩn sau: Hồi quy Box $(x,y,w,h)$, đánh giá độ trùng khớp \textbf{IoU}, và thuật toán lọc nhiễu \textbf{NMS}.
    \end{itemize}
\end{frame}

\end{document}
